\begin{frame}{Comparação: BCD, Gray e ASCII}
    \small

    \begin{center}
        \begin{tabular}{p{2cm}|p{3.2cm}|p{4.5cm}}
            \textbf{Código} &
            \textbf{Objetivo} &
            \textbf{Aplicações} \\ \hline

            BCD &
            Representar algarismos decimais &
            Displays, relógios, contadores e instrumentos \\[0.3cm]

            Gray &
            Reduzir alterações simultâneas de bits &
            Encoders, sensores e sistemas de posição \\[0.3cm]

            ASCII &
            Representar caracteres &
            Comunicação serial, computadores e textos
        \end{tabular}
    \end{center}

    \vspace{0.3cm}

    \begin{alertblock}{Resumo}
        \begin{itemize}
            \item BCD representa \textbf{dígitos};
            \item Gray representa \textbf{estados com transições seguras};
            \item ASCII representa \textbf{caracteres}.
        \end{itemize}
    \end{alertblock}

\end{frame}
